\section{Pengujian Dasar (Konsep Unit Test)}

\textbf{Unit test} memeriksa unit kode terkecil (fungsi, modul) secara terisolasi untuk memastikan perilaku yang benar. Test bersifat otomatis, repeatable, dan memberikan confidence saat refactoring: jika test tetap hijau, perubahan tidak merusak fungsionalitas. Framework test JavaScript populer: \textbf{Jest}, \textbf{Vitest}, dan \textbf{Mocha} \cite{mdnToolsTesting}.

\begin{webcode}[language=JavaScript, caption={Contoh unit test sederhana (konsep Jest)}]
// fungsi yang diuji
function hitungDiskon(harga, persen) {
  if (persen < 0 || persen > 100) throw new Error("Invalid");
  return harga - (harga * persen / 100);
}

// test cases (pola AAA: Arrange-Act-Assert)
test("diskon 20%", () => {
  expect(hitungDiskon(100000, 20)).toBe(80000);
});
test("diskon 0%", () => {
  expect(hitungDiskon(100000, 0)).toBe(100000);
});
test("persen negatif throws", () => {
  expect(() => hitungDiskon(100000, -5)).toThrow();
});
\end{webcode}

\begin{catatan}
\textbf{Test Coverage Bukan Segalanya.} 100\% coverage bukan jaminan kode bebas bug---test bisa saja tidak menguji edge case yang tepat. Fokus pada test yang bermakna: happy path (kasus normal), edge cases (batas, kosong, null), dan error cases. Test yang baik juga berfungsi sebagai dokumentasi: membaca test menjelaskan apa yang dilakukan fungsi \cite{mdnToolsTesting}.
\end{catatan}

